\documentstyle[%


righttag%


,11pt%


,amssymb%


,russian%               remove in Eng.


,izvru%                 Set izven% in Eng.


]{amsart}





%       Math definitions





\newcommand{\dg}[1]{{\bold #1}}


\newcommand{\re}{$d$-р.\,п.}


\newcommand{\mer}{\leqslant}


\newcommand{\bor}{\geqslant}


\newcommand{\rst}{\upharpoonright}


\newcommand{\op}{\oplus}


\newcommand{\dsum}{(D_1\op D_2)}


\newcommand{\draz}{(D_1-\nolinebreak D_2)}


\newcommand{\req}[2]{{\cal #1}_#2}


\newcommand{\con}[2]{#1\,\widehat{\ }\,#2}


\newcommand{\f}{\operatorname{first}}


\newcommand{\tre}{\operatorname{Req}}


\newcommand{\num}{\operatorname{num}}


\newcommand{\piv}[2]{$\langle\dg{#1},\dg{#2}\rangle$}


\newcommand{\tts}[1]{}





\theoremstyle{plain}


\theorembodyfont{\it}


\newtheorem{thmnonum}{\hskip\parindent Теорема}


\newtheorem{lems}{\hskip\parindent Лемма}[section]


\renewcommand{\thethmnonum}{}





%\hoffset=-15mm%                  For Eng.


 


\setcounter{page}{18}


\god{1998}


\nomer{7~(434)} %                        Remove in Eng.


%\nomer{Vol.\,42, No.\,7}%               Set in Eng.


%\str{\pageref{begin}--\pageref{end}}%  Set in Eng.


\udk{510.5 }


\author{\it А.А.\,Ефремов}


% NB:   ^^ remove in Eng.


\title{Изолированные сверху $d$-р.\,п. степени, II}





\thanks{Работа выполнена при частичной финансовой


поддержке Российского фонда фундаментальных исследований 


(номера проектов 93-011-16004, 96-01-00830) и гранта Новосибирского


Университета.}


\date{}


%\pagestyle{myheadings}%         Set in Eng.





\pagestyle{plain} %              Remove in Eng.


\begin{document}





%\thispagestyle{plain}%          Set in Eng.





\maketitle


\label{begin}





\section{Введение}





Множество $A\subseteq\omega$ называется $d$-{\it рекурсивно перечислимым}


(\re), если существуют рекурсивно перечислимые (р.п.) множества $A_1$ и


$A_2$ такие, что $A=A_1-A_2$. Тьюрингова степень называется \re{} степенью,


если она содержит \re{} множество; \re{} степень называется


{\it собственно} \re, 


если она не является р.\,п. степенью (не содержит р.\,п. множеств).





В статье продолжается изучение изолированных сверху (ИСВ) \re{} степеней.


Понятие ИСВ степени, введенное в первой части статьи~\cite{Efr}, является


естественным расширением понятия изолированной \re{} степени~\cite{CoYi}


(которому


в нашей терминологии соответствует понятие изолированности снизу (ИСН)).


Напомним, что \re{} степень $\dg{d}$ называется изолированной сверху, если


существует р.\,п. степень $\dg{b}>\dg{d}$ такая, что


между $\dg{d}$ и $\dg{b}$ нет р.\,п. степеней. В противном случае она


называется неизолированной сверху. В случае ИСВ говорим, что степени


$\dg{d}$ и $\dg{b}$


образуют изолированную сверху пару, и обозначаем \piv{d}{b}.





В первой части статьи изучался вопрос о плотности ИСВ степеней в структуре


р.\,п. степеней. Было установлено, что их ``достаточно много'': из основного


результата вытекало, что каждый класс скачков содержит бесконечно


возрастающую последовательность


ИСВ \re степеней.


В этой статье мы отвечаем полностью на вопрос о плотности, показывая, что


существует нерекурсивная р.\,п. степень, под которой {\it все} \re{} степени


являются неизолированными сверху, т.\,е. ИСВ степени в структуре р.\,п.


степеней не плотны.





\section{Обозначения и терминология}





Наши обозначения стандартны, и в основном мы следуем~\cite{So}. Используемая


терминология подробно описана в первой части статьи~\cite{Efr}, здесь


же напомним лишь некоторые определения.





При описании стратегий употребляем следующие термины: при выборе представителя


цикла слова ``достаточно большое число'' означают первое число, большее всех


упоминаемых в конструкции к данному моменту; {\it начинаем} цикл, позволяя ему


выполнить~(1) и перейти в~(2); {\it останавливаем} цикл, прекращая его работу,


определяя его запрет равным нулю и заставляя перейти в~(1); {\it разрушаем}


цикл, останавливая его и считая, что его часть функционала становится 


неопределенной; {\it инициализируем} стратегию, разрушая все ее циклы и


начиная


цикл 0; при инициализации все созданные связи и прикрепления уничтожаются;


цикл {\it работает}, переходя из некоторого (кроме (1)) состояния в


другое; стратегия {\it работает}, позволяя циклу с наименьшим номером, 


который может это сделать, работать (в противном случае ничего не делает).








\section{Основной результат}





Вначале мы доказываем следующий результат.


\begin{thmnonum}\label{DEN}


Существует нерекурсивное р.\,п. множество $A$ такое, что для любого \re{}


множества $\draz\mer_{\mathrm T}A$ либо найдется р.\,п. множество~$B$ такое,


что $\draz<_{\mathrm T}B<_{\mathrm T}\dsum$, либо $\draz\equiv_{\mathrm T}


\dsum$. 


\end{thmnonum}





{\bf Доказательство.} Пусть $\{V_e\}_{e\in\omega}$ --- некоторое эффективное


перечисление \re{} множеств. Заметим, что изолированная сверху пара \piv{d}{b}


характеризуется тем, что степени $\dg{d}$ и $\dg{b}$ имеют в структуре р.\,п.


степеней общие верхние конусы (это следует из теоремы Лахлана~\cite


{So}, с.\,164). Поэтому, чтобы построить р.\,п. множество $A$


с требуемыми свойствами,


достаточно удовлетворить для всех $e$ следующим требованиям:


\begin{itemize}


\item[$\req{P}{e}:$] $A\ne\Phi_e$, где $\{\Phi_e\}_{e\in\omega}$ ---


перечисление ч.\,р. функционалов;


\item[$\req{M}{e}:$] $V_e=\draz=\Phi_e^A\Longrightarrow


(\exists B_e)(\exists\Delta_e)\Bigl(\bigl(\draz=\Delta_e^{B_e}\bigr) \&\\


(\forall j)\bigl(B_e\ne\Phi_j^{\draz}\bigr)\&(\forall j)\bigl(\dsum\ne


\Phi_j^{B_e}\bigr)\Bigr)\vee(\exists\Gamma_e)\Bigl(\dsum=\Gamma_e^{\draz}


\Bigr)$, 


\end{itemize}





\noindent где $\{V_e,\Phi_e\}$ --- перечисление пар, состоящих из


\re{} множества $V_e$ и ч.\,р. функционала $\Phi_e$, $B_e$ --- р.\,п.


множество, $\Delta_e$, $\Gamma_e$ --- ч.\,р. функционалы.





Рассмотрим основные модули для удовлетворения этих требований.





\underline{\bf Основной модуль для $\req{P}{e}$}: $A\neq\Phi_e$





\noindent Этот вид требований обеспечивает нерекурсивность


множества $A$. Стратегия работает следующим образом.


\begin{itemize}


\item[(1)] Выбираем представителя цикла $x$ достаточно большим.


\item[(2)] Ждем шага $s: \Phi_e(x)\downarrow=0  [s]$.


\item[(3)] Перечисляем $x$ в $A$.


\item[(4)] Конец работы.


\end{itemize}





\underline{\bf Выходы $\req{P}{e}$-стратегии}





\begin{itemize}


\item[1)] Стратегия ждет всегда в~(2). Выход обозначим через {\it 1}.


\item[2)] Стратегия проходит через~(3) и остается навсегда в~(4).


Выход обозначим через {\it 0}.


\end{itemize}





$\req{P}{e}$-стратегия действует один раз (без учета инициализаций со стороны


других стратегий) и требует лишь того, чтобы запрет стратегий с большим


приоритетом на $A$\, либо был конечным, либо бесконечно много раз снимался.





Требование $\req{M}{e}$ очень сложно для удовлетворения в том виде, в


котором мы его записали. Поэтому мы его разбиваем на три части:


требование $\req{R}{e}$ и подтребования $\req{N}{{e,j}}$ и


$\req{S}{{e,j}}$.


\begin{itemize}


\item[$\req{R}{e}:$] $\draz=\Phi_e^A\Longrightarrow


(\exists B_e)(\exists\Delta_e)\Bigl(\draz=\Delta_e^{B_e}\Bigr) \bigvee


(\exists\Gamma_e)\Bigl(\dsum=\Gamma_e^{\draz}\Bigr)$.


\item[$\req{N}{{e,j}}:$] $\draz=\Phi_e^A\Longrightarrow


\dsum\neq\Phi_j^{B_e}$.


\item[$\req{S}{{e,j}}:$] $\draz=\Phi_e^A\Longrightarrow


B_e\neq\Phi_j^{\draz}$.


\end{itemize}





Введем вспомогательные о.\,р. функции:





$\ell(e,s)=\max\{x\mid(\forall y<x)\bigl[\draz(y)=\Phi_e^A(y) [s]\bigr]\}$,





$m(e,s)=\max\{\ell(e,t)\mid t\mer s\}$.








Будем говорить, что шаг $s$ является $e$-расширяющим, если $s=0$ или


$\ell(e,s)>m(e,s-1)$.





\underline{\bf Основной модуль для $\req{R}{e}$}





\noindent $\req{R}{e}:$ $\draz=\Phi_e^A\Longrightarrow


(\exists B_e)(\exists\Delta_e)\Bigl(\draz=\Delta_e^{B_e}\Bigr) \bigvee


(\exists\Gamma_e)\Bigl(\dsum=\Gamma_e^{\draz}\Bigr)$.





\noindent Стратегия для одиночного $\req{R}{e}$\, состоит в следующем:


если верна посылка, то мы хотим выполнить первый дизъюнкт, т.\,е. построить


всюду определенный функционал $\Delta_e$ такой, что $\draz=\Delta_e^{B_e}$.


Поэтому для каждого $x$ ждем первый шаг $s$ такой, чтобы $\ell(e,s)>x$, и


определяем $\Delta_e^{B_e}(x)=\draz(x)$ с достаточно большим $\delta_e(x)$.


При необходимости корректируем функционал $\Delta_e^{B_e}$ в точке


$x$, перечисляя $\delta_e(x)$ в $B_e$ (при этом делаем неопределенным


$\Delta_e^{B_e}(y)$ во всех $y>x$).





\underline{\bf Выходы $\req{R}{e}$-стратегии}





\begin{itemize}


\item[1)] Конечный выход нашей стратегии, соответствующий ситуации, когда


$\lim_s\ell(e,s)<\infty$. Требование $\req{R}{e}$ выполнено, причем нет


необходимости выполнять требования $\req{N}{{e,j}}$ и


$\req{S}{{e,j}}$, $j\in\omega$.


Выход стратегии обозначим в этом случае через {\it 1}.


\item[2)] Бесконечный выход нашей стратегии, соответствующий ситуации, когда


$\lim_s\ell(e,s)=\infty$.


Тогда функционал $\Delta_e^{B_e}$ является всюду определенным, и


ниже этого выхода возможны бесконечные нарушения на $B_e$.


Выход стратегии обозначим в этом случае через {\it 0}.


\end{itemize}





\underline{\bf Основной модуль для $\req{S}{{e,j}}$}





\noindent $\req{S}{{e,j}}:\draz=\Phi_e^A\Longrightarrow B_e\neq\Phi_j^{\draz}$





\noindent Требования $\req{S}{{e,j}}$ на дереве стратегий будут расположены


только под бесконечным выходом $\req{R}{e}$-стратегии. Поэтому будем


предполагать, что $\draz=\Phi_e^A$.





Стратегия работает следующим образом.


\begin{itemize}


\item[(1)] Выбираем представителя стратегии $x$ достаточно большим.


\item[(2)] Ждем шага $s:\Phi_j^{\draz}(x)=0\;\&\;\draz\rst(\varphi_j(x)+1)=


\Phi_e^A \rst(\varphi_j(x)+1) [s]$.


\item[(3)] Перечисляем $x$ в $B_e$, устанавливаем запрет $r[s+1]=


\varphi_e(\varphi_j(x)) [s]$ на~$A$.


\item[(4)] Конец работы.


\end{itemize}


 


\underline{\bf Выходы $\req{S}{{e,j}}$-стратегии}





\begin{itemize}


\item[1)] Стратегия всегда ждет в~(2). Работы никакой не производилось.


Выход стратегии в этом случае обозначим через {\it 1}.


\item[2)] Стратегия проходит через~(3) и остается навсегда в~(4).


Запрет на $A$ конечный. Выход обозначим через {\it 0}.


\end{itemize}


$\req{S}{{e,j}}$-стратегия действует один раз (без учета инициализаций со


стороны


других стратегий) и требует лишь того, чтобы запрет стратегий с большим


приоритетом на $B_e$ либо был конечным, либо бесконечно много раз снимался.





\underline{\bf Основной модуль для $\req{N}{{e,j}}$}





\noindent $\req{N}{{e,j}}:\draz=\Phi_e^A\Longrightarrow \dsum\ne\Phi_j^{B_e}$.





\noindent Требования $\req{N}{{e,j}}$ на дереве стратегий будут расположены


только под бесконечным выходом $\req{R}{e}$-стратегии. Поэтому будем


предполагать, что $\draz=\Phi_e^A$. Именно между $\req{N}{{e,j}}$-


и $\req{R}{e}$-стратегиями возникает первый конфликт нашей конструкции:


$\req{N}{{e,j}}$-стратегия будет порождать запреты на $B_e$, а


$\req{R}{e}$-стратегия может бесконечно нарушать их.


В этом случае $\req{N}{{e,j}}$-стратегия будет обходить эту трудность,


выполняя второй дизъюнкт


требования $\req{R}{e}$ (каждая $\req{N}{{e,j}}$-стратегия


строит свой вариант функционала


$\Gamma_e$).





%Вначале рассматриваем случай, когда $\req{N}{{e,j}}$-стратегия непосредственно


%следует за $\req{R}{e}$-стра\-те\-ги\-ей.





$\req{N}{{e,j}}$-стратегия работает по циклам, каждый из которых


работает следующим образом:


\begin{itemize}


\item[] \underline{Цикл $k$}


\item[(1)] Выбираем представителя цикла $x_k=2n_k+1$ достаточно большим.


\item[(2)] Ждем шага $s: \dsum\rst(x_k+1)=\Phi_j^{B_e}\rst(x_k+1)\&


\draz\rst(y_k+1)=\Phi_e^A\rst(y_k+1) [s]$, где


$$


y_k=\max\{n_k,\max\{z\mid\delta_e(z)<\varphi_j(x_k)[s]\}\}.


$$


\item[(3)] Определяем функционал


$\Gamma_e^{\draz}\rst(x_k+1)=\dsum\rst(x_k+1)[s+1]$


с $\gamma_e(x_k)[s+1]=y_k$ в точках, в которых он еще


не определен.


Определяем запреты $r_1[s+1]=\varphi_j(x_k)[s]$ на $B_e$ и


$r_2(k)[s+1]=\varphi_e(y_k)[s]$ на $A$.


Начинаем $(k+1)$-й цикл. Одновременно продолжаем работу нашего цикла,


переходя в~(4).


\item[(4)] Ждем шага $t>s:\dsum\rst(x_k+1)[t]\neq\dsum\rst(x_k+1)[s]$, но


$\draz\rst(y_k+1)[t]=\draz\rst(y_k+1)[s]$.


Пока такого шага $t$ нет, мы, следуя изменениям $\dsum\rst(x_k+1)$,


имеем возможность корректировать функционал $\Gamma_e^{\draz}$ в


соответствующих точках, т.\,к. $\draz\rst(y_k+1)$ тоже меняется.


Как только такой шаг $t$ нашелся, останавливаем циклы с номерами $>k$ и


идем в~(5).


\item[(5)] Конец работы.


\end{itemize}





\underline{\bf Выходы $\req{N}{{e,j}}$-стратегии}





\begin{itemize}


\item[1)] Существует шаг $s_0$, после которого ни один цикл не работает.


Тогда некоторый цикл $k_0$ остается навсегда в~(2) или в~(5).


Требование $\req{N}{{e,j}}$ удовлетворено. Общие запреты всех циклов


на $B_e$ и $A$ имеют конечный предел. Заметим, что запрет на


$B_e$ не мешает $\req{R}{e}$-стратегии корректировать свой


функционал, т.\,к. значения $\draz$ для всех $y:\delta_e(y)<\varphi_j(x_k)$


являются такими же, какими были при определении


$\Delta_e^{B_e}(y)$, а дальнейшие изменения $\draz\rst(y+1)$


контролируются запретом $r_2$ на $A$.


Обозначим этот выход через {\it 1} (конечный) (для более простого


вида дерева стратегий объединяем два конечных выхода в одном).


\item[2)] Работает бесконечно много циклов, каждый из которых ждет в~(4).


Общий запрет всех циклов на $B_e$ и на $A$ может иметь бесконечный


предел. Требование $\req{N}{{e,j}}$ не выполнено, но построен


всюду определенный функционал $\Gamma_e^{\draz}=\dsum$,


тем самым $\req{N}{{e,j}}$-стратегия полностью удовлетворяет


требование $\req{R}{e}$ и ниже этого выхода не


нужно больше ни $\req{S}{{e,i}}$, ни $\req{N}{{e,i}}$, $i>j$.


Выход бесконечный, обозначим его через {\it 0}.


\end{itemize}





Теперь ниже бесконечного выхода $\req{N}{{e,j}}$-стратегии возникает


вторая проблема: т.\,к. запрет на $A$ может быть бесконечным, то


могут не выполниться ${\cal P}$-стратегии.





\underline{\bf $\cal P$-стратегия ниже бесконечного выхода


$\cal N$-стратегии}





Фиксируем произвольную $\cal N$-стратегию и некоторую $\cal P$-стратегию,


расположенную под бесконечным выходом нашей стратегии.


Обозначим через $x$ текущий представитель $\cal P$-стратегии, и пусть на


шаге $s$ цикл $k$\; $\cal N$-стратегии установил запрет на $A$\; $r_2[s]>x$.


Предположим далее, что на шаге $t\bor s$\; $\cal P$-стратегия хочет


перечислить~$x$ в~$A$, чтобы удовлетворить $\req{P}{e}$. Основная идея


решения этой проблемы состоит в том, что мы будем разрешать


$\cal P$-стратегии перечислять элемент~$x$ в~$A$ в этом случае, соблюдая


одно условие: каждый цикл $\cal N$-стратегии может нарушаться нижними


$\cal P$-стратегиями не более одного раза. Перечислив $x$ в~$A$, мы


нарушаем запрет $r_2(k)$ и тем самым, возможно, разрушаем некоторое


вычисление $\Phi_e^A\rst(y_k+1) [s]$. Цикл~$k$ должен обработать


эту ситуацию, поэтому изменим цикл $\cal N$-стратегии следующим


образом:





\begin{itemize}


\item[(1),] (2) и (3) остаются без изменений.


\item[(4)$\phantom{,}$] Ждем шаг $t>s$ такой, что выполняется либо а),


либо б). Пока такого шага $t$ нет, мы, следуя изменениям $\dsum\rst(x_k+1)$,


имеем возможность корректировать функционал $\Gamma_e^{\draz}$ в


соответствующих точках, т.\,к. $\draz\rst(y_k+1)$ тоже меняется.


\begin{itemize}


\item[а)]$\dsum\rst(x_k+1)[t]\neq\dsum\rst(x_k+1)[s]$,но


$\draz\rst(y_k+1)[t]=\draz\rst(y_k+1)[s]$. Тогда останавливаем циклы


с номерами $>k$ и идем в~(5).


\item[б)] Изменилось $A\rst[r_2(k-1),r_2(k)]$ и $\draz\rst[y_{k-1}+1,y_k]$.


Тогда разрушаем циклы с номерами $>k$ и начинаем $(k+1)$-й


цикл. Переходим в~(5).


\end{itemize}


\item[(5)] Конец работы.


\end{itemize}





Разница между пп.\,(4а) и (4б) состоит в следующем.


Пока цикл ждет в~(4), сохраняются две возможности удовлетворения


требований $\req{R}{e}$ и $\req{N}{{e,j}}$: либо удовлетворяем


$\req{N}{{e,j}}$ без корректировки функционала $\Delta_e$ требования


$\req{R}{e}$, либо не удовлетворяем требования $\req{N}{{e,j}}$,


но строим функционал $\Gamma_{e,j}$, тем самым полностью удовлетворяя


требование $\req{R}{e}$. Проходя в~(5) через~(4а), реализуем первую


возможность; проходя в~(5) через~(4б), делаем шаг по реализации второй


возможности.





Рассмотрим подробнее п.\,(4б). Если цикл~$k$ попал в~(4б), это означает,


что нижняя $\cal P$-стратегия перечислила элемент~$x$ в~$A$, который


повлиял на вычисления, т.\,к. равенство на отрезке $[y_{k-1}+1,y_k]$


восстановилось только после изменения $\draz$. Тем самым имеем


возможность корректировать функционал $\Gamma_{e,j}$ в соответствующих


точках. Если же в будущем $\draz$ на отрезке $[y_{k-1}+1,y_k]$ вернется


к предыдущему значению, просто добиваемся неравенства $\draz\neq\Phi_e^A$,


т.\,к. контролируем $A\rst\varphi_e(y_k)$.





Выходы стратегии остаются прежними, внесем только небольшие изменения.


\begin{itemize}


\begin{itemize}


\item[\underline{Выход 1}.] Существует шаг $s_0$, после которого ни


один цикл не работает. Тогда некоторый цикл $k_0$ остается навсегда


в~(2) или в~(5), пройдя через~(4а).


\item[\underline{Выход 2}.] Работает бесконечно много циклов, каждый из


которых либо ждет в~(4), либо находится в~(5), пройдя через~(4б).


\end{itemize}


\end{itemize}





Случай, когда $\cal P$-стратегия находится под бесконечным выходом


нескольких $\cal N$-стратегий, никаких новых проблем не ставит. Подчеркнем,


что вышеописанная стратегия взаимодействия $\cal P$- и $\cal N$-стратегий


будет работать при условии, что каждый цикл $\cal N$-стратегии нарушается


нижними $\cal P$-стратегиями не более одного раза. Реализация этой идеи


будет описана в Полной конструкции.





\underline{\bf Дерево стратегий}





Пусть $\Lambda=\{\text{\it 0}<_{\Lambda}\text{\it 1}\}$ ---


множество возможных выходов


наших стратегий. Тогда деревом стратегий будем называть множество


$T=\Lambda^{<\infty}$. Следующее, что должны сделать, это назначить


каждой вершине $\alpha\in T$ некоторое требование.


Назначение требований вершинам дерева будем проводить индукцией


по $n=|\alpha|$, определяя функцию $\tre(\alpha)$. При этом будем


использовать вспомогательную функцию


$L$, которая назначает каждой вершине некоторый упорядоченный список 


требований, которые ждут прикрепления к вершинам дерева. Вначале этот список


будет включать все стратегии, упорядоченные в естественном порядке. По мере 


продвижения по дереву $T$ будем добавлять или убирать некоторые 


стратегии. Обозначим через $L_0$ список всех стратегий в естественном


порядке


$$


L_0=\{\req{P}{0},\req{R}{0},\req{S}{{0,0}},\req{N}{{0,0}},\req{P}{1},


\req{R}{1},\req{S}{{0,1}},\req{N}{{0,1}},\dotsc,\req{P}{e},\req{R}{e},


\req{S}{{l(e),r(e)}},\req{N}{{l(e),r(e)}},\dotsc \};


$$ 


через $\f (L)$ --- первый элемент в списке $L$; через $\num (X)$


--- номер требования $X$.





Теперь опишем стратегию прикрепления требований к вершинам.


 


Пусть $n=|\alpha|=0$. Определяем $L(\alpha)=L(\lambda)=L_0,


\tre (\lambda)=\f (L(\lambda))$, т.\,е. вершине $\lambda$ назначаем требование


$\req{P}{0}$.





Пусть $n=|\alpha|>0$. Тогда $\alpha=\con{\beta}{a}$ для некоторого


$\beta$, причем $L(\beta)$ и $\tre (\beta)$ определены.


Определяем вначале $L(\alpha)$. Возможны следующие случаи.


\begin{itemize}


\item[(1)] $\tre(\beta)=\req{P}{e}$ для некоторого $e\in\omega$, 


$a=0,1$. Полагаем $L(\alpha)=L(\beta)\setminus\{\tre(\beta)\}$.


\item[(2)] $\tre(\beta)=\req{S}{{e,j}}$ для некоторых $e,j\in\omega$,


$a=0,1$. Полагаем $L(\alpha)=L(\beta)\setminus\{\tre(\beta)\}$.


\item[(3)] $\tre(\beta)=\req{R}{e}$ для некоторого $e\in\omega$.


\begin{itemize}


\item[(a)] $a=0$. Полагаем $L(\alpha)=L(\beta)\setminus\{\tre(\beta)\}$.


\item[(б)] $a=1$. Полагаем $L(\alpha)=L(\beta)\setminus\bigl(\{\tre(\beta)\}


\bigcup\{X\mid X\in L(\beta)\&{} \num(X)=(e,j)$


для некоторого $j\in\omega\}\bigr)$.


\end{itemize}


\item[(4)] $\tre(\beta)=\req{N}{{e,j}}$ для некоторых $e,j\in\omega$.


\begin{itemize}


\item[(a)] $a=0$. Полагаем $L(\alpha)=L(\beta)\setminus\bigl(\{\tre(\beta)\}


\bigcup\{X\mid X\in L(\beta)\&{} \num(X)=(e,j)$


для некоторого $j\in\omega\}\bigr)$.


\item[(б)] $a=1$. Полагаем $L(\alpha)=L(\beta)\setminus


\{\tre(\beta)\}$.


\end{itemize}


\end{itemize}





Окончательно определяем $\tre(\alpha)=\f(L(\alpha))$.





\noindent На этом описание стратегии прикрепления требований


к вершинам закончено.


 


Теперь $\alpha$-стратегией будем называть вариант основного модуля для 


требования, которое назначено вершине $\alpha$, с некоторыми изменениями,


описанными в конструкции; $S$-стратегией, где


$S\in\{{\cal P},{\cal R},{\cal N},{\cal S}\}$ будем называть стратегию


$\alpha$, если она работает над требованием вида $\req{P}{e}$, 


$\req{R}{e}$, $\req{N}{{e,j}}$, $\req{S}{{e,j}}$ соответственно. Отметим, что


каждая $\req{R}{e}$-стратегия $\alpha$ будет строить свои варианты р.\,п.


множества $B_e$ и функционала $\Delta_e$. Думаем, что никакой путаницы не


возникнет, и поэтому не будем добавлять лишних индексов к нашим множествам.


Шаг $s$ назовем


$\alpha$-шагом, если стратегия $\alpha$ имела возможность работать на


шаге $s$. Для работы со стратегиями $\alpha$ на дереве стратегий


изменим определения функций длины $\ell(e,s)$, $m(e,s)$


следующим образом. Пусть $\tre(\alpha)=\req{R}{e}$. Тогда, если $s$ ---


$\alpha$-шаг, то $\ell(\alpha,s)=\ell(e,s)$; в противном случае


$\ell(\alpha,s)=\ell(e,t)$, где $t$ --- наибольший $\alpha$-шаг,


меньший $s$. Аналогично определяем $m(\alpha,s)$.


\pagebreak





\underline{\bf Полная конструкция}


\medskip





\noindent\underline{\bf Шаг 0}. $A_0=\emptyset$,


для всех $e$ все $B_{e,0}=\emptyset$,


$\Delta_e$, $\Gamma_{e,j}$ не определены, все стратегии $\alpha\in T$


инициализированы.





\noindent\underline{\bf Шаг $s+1$}. Работаем по подшагам $t\mer s$,


причем на каждом подшаге $t$ одна из стратегий $\alpha$ длины


$t$ будет иметь возможность работать (при $t=0$ работает стратегия


$\lambda$, которой назначено требование $\req{P}{0}$). Возможны четыре


случая.


 


$\tre(\alpha)=\req{P}{e}$. Стратегия $\alpha$ работает над требованием


${\cal P}_e$, как описано в основном модуле для ${\cal P}_e$.


После этого определяем выход  $o$ стратегии $\alpha$ следующим образом:


если стратегия находится в~(4), то выходом будет {\it 0}, в противном


случае --- {\it 1}.


Инициализируем стратегии $\beta>_{\mathrm L}\con{\alpha}{o}$.


Далее, если стратегия только что прошла через~(3) в~(4) (перечислила некоторый


элемент в $A$), то определяем $\delta_{s+1}=\alpha$ и переходим к следующему


шагу $s+2$. В противном случае определяем следующую стратегию $\beta$,


которая имеет возможность работать на подшаге $t+1 : \beta=\con{\alpha}{o}$.


Переходим к подшагу $t+1$.





$\tre(\alpha)=\req{R}{e}$. Стратегия $\alpha$ работает над требованием


${\cal R}_e$. Работу производим только в том случае, если шаг $s$ является


$e$-расширяющим. В первую очередь проверяем, есть ли точки, в которых


необходимо корректировать функционал $\Delta_e$.


Если есть, то выбираем такую наименьшую точку $x$ и перечисляем


$\delta_e(x)$ в $B_e$. После этого функционал $\Delta_e$ во всех


точках $y>x$ считается неопределенным.


Определяем $\delta_{s+1}=\alpha$ и переходим к следующему шагу $s+2$.


Если корректировать нечего, то стратегия определяет функционал $\Delta_e$


в точках $x$, для которых $\ell(\alpha,s)>x$ и в которых он не определен.


После этого определяем выход $o$ стратегии $\alpha$ следующим образом:


если шаг $s$ был $e$-расширяющим, то выходом считаем {\it 0}, в


противном случае --- {\it 1}.


Определяем следующую стратегию $\beta$, которая имеет возможность работать на


подшаге $t+1: \beta=\con{\alpha}{o}$. Переходим к подшагу $t+1$.


  


$\tre(\alpha)=\req{S}{{e,j}}$. Стратегия $\alpha$ работает над требованием


${\cal S}_{e,j}$, как описано в основном модуле для  ${\cal S}_{e,j}$.


После этого определяем выход  $o$ стратегии $\alpha$ следующим образом:


если стратегия находится в~(4), то выходом будет {\it 0}, в противном


случае --- {\it 1}.


Инициализируем стратегии $\beta>_{\mathrm L}\con{\alpha}{o}$.


Далее, если стратегия только что прошла через~(3) в~(4) (перечислила некоторый


элемент в $B_e$), то определяем $\delta_{s+1}=\alpha$ и переходим к


следующему шагу


$s+2$. В противном случае определяем следующую стратегию $\beta$,


которая имеет возможность работать на подшаге $t+1 : \beta=\con{\alpha}{o}$.


Переходим к подшагу $t+1$.





$\tre(\alpha)=\req{N}{{e,j}}$. Стратегия $\alpha$ работает над требованием


${\cal N}_{e,j}$, как описано в основном модуле для ${\cal N}_{e,j}$.


Если работал цикл $k$ и установил новый запрет $r_2(k)$, то


ищем $\cal P$-стратегии $\beta\supseteq\con{\alpha}{{\it 0\/}}$ такие, что


их представители $x(\beta)<r_2(k)$ и $x(\beta)\not\in A$. Если таких нет,


ничего не делаем. Если есть,


выбираем $<$-наименьшую такую $\beta_0$. Связываем $\beta_0$ с циклом $k$ и


инициализируем $\cal P$-стратегии $\gamma>\beta_0$ такие, что


$\gamma\supseteq\con{\alpha}{\text{\it 0}}$.


После этого определяем выход $o$ стратегии $\alpha$ следующим образом:


если стратегия начала новый цикл, то выходом будет


{\it 0}, в противном случае --- {\it 1}.


Инициализируем стратегии $\beta>_{\mathrm L}\con{\alpha}{o}$.


Далее, определяем следующую стратегию $\beta$, которая имеет возможность


работать на подшаге $t+1 : \beta=\con{\alpha}{o}$.


Переходим к подшагу $t+1$.





В каждом из рассмотренных случаев при $t=s$ следующую стратегию не выбираем,


а, завершая шаг $s+1$, определяем $\delta_{s+1}=\alpha$ и переходим к


следующему шагу $s+2$.





\noindent Конструкция закончена.





\medskip


\underline{\bf Проверка конструкции}


\medskip


   


\noindent Для завершения доказательства теоремы докажем ряд технических лемм.


\addtocounter{section}{-2}


 


Определим истинный путь $f\in [T]$ как самую левую бесконечную ветвь


дерева $T$, посещаемую $\delta_s$ в течение конструкции бесконечное


число раз, т.\,е. для всех $n$, если $\alpha=f\rst n$, то $f(n)=a$


означает окончательный выход стратегии $\alpha$. Покажем, что истинный


путь в нашей конструкции определяется корректно. 


 


\begin{lems}


Истинный путь $f$ существует.


\end{lems}





\begin{pf}


Во-первых, очевидно, что $\lim_s|\delta_s|=\infty$. Во-вторых, наше


дерево стратегий бинарно. Поэтому ясно, что истинный путь $f$


определяется однозначно.


\end{pf}





\begin{lems}\label{INI}


Все вершины $\alpha\subset f$ инициализируются конечное число раз.


\end{lems}





\begin{pf}


Пусть $\alpha$ --- не $\cal P$-стратегия.


Стратегию $\alpha\subset f$ могут инициализировать только такие стратегии


$\beta$, для которых либо $\beta<_{\mathrm L}\alpha$,


либо $\beta\subset\alpha$.


В первом случае, т.\,к. $\alpha\subset f$, то, по определению


истинного пути, $\beta$ может это делать только конечное число раз.


Во-втором случае по нашей конструкции никакая $\beta$ не может


инициализировать $\alpha\supset\beta$ вообще.


Пусть теперь $\alpha$ --- $\cal P$-стратегия. По определению истинного пути


существует шаг $s_0$ такой, что для $t\bor s_0$\;


$\delta_t\not<_{\mathrm L}\alpha$.


Учитывая это и тот факт, что $\cal P$-стратегий $\beta\subset\alpha$ конечное


число, заключаем, что существует шаг $s_1\bor s_0$, к которому все возможные


прикрепления стратегий $\beta<\alpha$ к циклам $\cal N$-стратегий уже сделаны.


После шага $s_1$ стратегия $\alpha$ больше инициализироваться не будет.


\end{pf}





Введем некоторые обозначения. Пусть $\tre(\alpha)=\req{N}{{e,j}}$. Обозначим


через $R_1(\alpha,s)$ запрет стратегии $\alpha$ на $B_e$ после шага $s$;


через $R_2(\alpha,s)$ --- запрет стратегии $\alpha$ на $A$ после шага $s$.


Для $\cal S$-стратегий $\beta$\; $R_2(\beta,s)$ определяем аналогично;


для всех $s$ полагаем


$R_1(\beta,s)=0$.


Для $\cal R$-стратегий и $\cal P$-стратегий $\gamma$


для всех $s$\; $R_1(\gamma,s)=R_2(\gamma,s)=0$.





\begin{lems}


Требования $\req{P}{e}$ для всех $e$ вдоль истинного пути удовлетворяются.


\end{lems}





\begin{pf}


Стратегия прикрепления требований к вершинам дерева стратегий гарантирует,


что вдоль любой бесконечной ветви расположены все требования $\req{P}{e}$.





Рассмотрим любую $\cal P$-стратегию $\alpha\subset f$.


По лемме~\ref{INI} существует


шаг $s$, после которого $\alpha$ уже не инициализируется.


Пусть


$\con{\beta_0}{\text{\it 0}}\subseteq\con{\beta_1}{\text{\it 0}}


\subseteq\dotsb\subseteq\con{\beta_n}{\text{\it 0}}\subseteq\alpha\subset f$,


где $\beta_i$ --- все такие $\cal N$-стратегии, что $\alpha$ находится под


их бесконечными выходами.


После шага $s$ все прикрепления стратегии $\alpha$ к некоторым циклам $k_i$


стратегий $\beta_i$ будут окончательными. Поэтому существует шаг $t\bor s$,


к которому все возможные связи стратегии $\alpha$ будут определены.


Таким образом, после шага $t$ на $\alpha$-шагах стратегии $\alpha$ уже


ничто не будет мешать выполнить требование $\tre(\alpha)$.


\end{pf}





\begin{lems}


Требования $\req{M}{e}$ для всех $e$ вдоль истинного пути удовлетворяются.


\end{lems}





\begin{pf}


Фиксируем требование $\req{M}{e}$. Ему соответствуют $\req{R}{e}$-,


$\req{S}{{e,j}}$- и $\req{N}{{e,j}}$-стратегии, $j\in\omega$.


Стратегия прикрепления требований вершинам дерева стратегий гарантирует, что


вдоль любой бесконечной ветви расположены все требования $\req{R}{e}$.





Пусть $\alpha\subset f$ такова, что $\tre(\alpha)=\req{R}{e}$.





Если $\con{\alpha}{\text{\it 1}}\subset f$, то


$\lim_s \ell(\alpha,s)<\infty$, и


требование $\req{R}{e}$ выполнено. Более того, ниже конечного выхода


$\req{R}{e}$ не нужны ни $\req{S}{{e,j}}$-, ни $\req{N}{{e,j}}$-стратегии.


Таким образом, требование $\req{M}{e}$ в этом случае выполнено.





Пусть $\con{\alpha}{\text{\it 0}}\subset f$.


Вначале покажем, что в этом случае


$\Delta_e^{B_e}$ будет всюду определенным и правильно вычислять $\draz$.


Действительно: в данном случае существует бесконечно много $e$-расширяющих


шагов, для любого $x$ $\draz\rst x$ может меняться лишь конечное число раз,


поэтому по конструкции $\Delta_e^{B_e}$ будет всюду определенным. Так как на


каждом $e$-расширяющем шаге мы в первую очередь корректируем $\Delta_e^{B_e}$,


то легко видеть, что он правильно вычисляет $\draz$ во всех точках (начиная с


некоторого шага $s$, после которого $\alpha$ больше не инициализируется).





Далее, возможны два случая:


1) существует $\req{N}{{e,j}}$-стратегия $\beta_0$ такая, что


$\con{\beta_0}{\text{\it 0}}\subset f$;


2) для всех $\req{N}{{e,j}}$-стратегий $\beta$\;


$\con{\beta}{\text{\it 1}}\subset f$.





1) По лемме~\ref{INI} существует шаг $s$, после которого $\beta_0$


не инициализируется.


Так как $\con{\beta_0}{\text{\it 0}}\subset f$, то работает бесконечно


много циклов стратегии $\beta_0$, каждый из которых либо ждет всегда в~(4),


либо проходит в~(5) через~(4б). Это означает, что начиная с шага $s$,


строящийся функционал $\Gamma_{e,j}^{\draz}$ будет всюду определенным


и правильно вычислять $\dsum$. Тем самым $\beta_0$ удовлетворяет требование


$\req{R}{e}$, выполняя его второй дизъюнкт. Ниже бесконечного выхода


$\beta_0$ уже нет ни $\req{S}{{e,i}}$-, ни $\req{N}{{e,i}}$-стратегий,


$i>j$. Требование $\req{M}{e}$ выполнено.





2) Фиксируем любую стратегию $\beta\supseteq\con{\alpha}{\text{\it 0}}$ такую,


что $\tre(\beta)=\req{N}{{e,j}}$. По лемме~\ref{INI} существует шаг $s$,


после которого $\beta$ не инициализируется. Так как


$\con{\beta}{\text{\it 1}}\subset f$, то некоторый цикл стратегии $\beta$


после шага $s$ либо


всегда ждет в~(2), либо прошел в~(5) через~(4а). Таким образом, требование


$\req{N}{{e,j}}$ выполнено. Отметим, что в этом случае существуют конечные


пределы запретов на $B_e$ и $A$: $R_1(\beta)=\lim_s R_1(\beta,s)$ и


$R_2(\beta)=\lim_s R_2(\beta,s)$. Таким образом, все требования


$\req{N}{{e,j}}$ в этом случае выполняются, порождая


конечные запреты ниже своего конечного выхода. Напомним, что при этом


$\cal N$-стратегия не мешает $\req{R}{e}$-стратегии $\alpha$ строить


всюду определенный


функционал $\Delta_e^{B_e}$, правильно вычисляющий $\draz$.


Теперь докажем, что и все требования


$\req{S}{{e,j}}$, лежащие на истинном пути, удовлетворяются в этом случае.


Фиксируем любую стратегию $\gamma\supseteq\con{\alpha}{\text{\it 0}}$ такую,


что $\tre(\gamma)=\req{S}{{e,j}}$. По лемме~\ref{INI} существует шаг $s$,


после которого $\gamma$ не инициализируется. Рассмотрим следующий запрет


на~$B_e$:


$$


R=\max\{R_1(\beta)\mid \con{\alpha}{\text{\it 0}}\subseteq\beta<\gamma


\ \;\&\ \;


\num(\tre(\beta))=(e,i)\text{ для некоторого }i\in\omega\}.


$$


Из вышеизложенных фактов следует, что он конечен и к шагу $s$ уже будет


достигнут. Следовательно, после шага $s$ стратегии $\gamma$ уже ничто не будет


мешать выполнить требование $\req{S}{{e,j}}$.





Таким образом, $\req{R}{e}$-стратегия построила всюду определенный функционал


$\Delta_e^{B_e}=\draz$, все $\req{S}{{e,j}}$ и $\req{N}{{e,j}}$ выполнены.


Следовательно, $\req{M}{e}$ выполнено.


\end{pf}





Доказательство теоремы закончено.





Теперь можем получить наш основной результат. Вначале напомним, что в первой


части статьи~\cite{Efr} был доказан критерий ИСВ пары:


собственно \re{} степень $\dg{d}$ и р.\,п. степень $\dg{b}>\dg{d}$ образуют


изолированную сверху пару тогда и только тогда, когда существует \re{}


множество $\draz\in\dg{d}$ такое, что $\dsum\in\dg{b}$ и между


степенями $\text{deg\,}\draz$ и $\text{deg\,}\dsum$ нет р.\,п. степеней.


Из этого критерия и нашей теоремы легко получаем следующее





\begin{cor}\label{NEI}


Существует нерекурсивная р.\,п. степень $\dg{a}$ такая, что все \re{} степени


$\dg{d}<\dg{a}$ являются неизолированными сверху.


\end{cor}





\begin{pf}


Пусть $\dg{a}=\text{deg}(A)$, где $A$ --- построенное нами


в теореме нерекурсивное


р.\,п. множество. Пусть $\dg{d}<\dg{a}$ --- собственно \re{} степень.


По критерию ИСВ пары, если бы $\dg{d}$ была ИСВ степенью,


тогда существовало бы


\re{} множество $\draz\in\dg{d}$ такое, что между $\draz$ и $\dsum$ нет р.\,п.


множеств


(по тьюринговой сводимости). Но по теореме для {\it любого} \re{}


множества $\draz<_{\mathrm T}A$ либо найдется р.\,п. множество


$B$: $\draz<_{\mathrm T}


B<_{\mathrm T}\dsum$, либо $\draz\equiv_{\mathrm T}\dsum$.


Так как степень $\dg{d}$


является собственно \re{} степенью, то должно выполняться первое свойство. 


Следовательно, степень $\dg{d}$ неизолирована сверху.


\end{pf}





\begin{cor}


ИСВ степени в структуре р.\,п. степеней не плотны.


\end{cor}





\begin{cor}


Существует неизолированная сверху собственно \re{} степень.


\end{cor}





\begin{thebibliography}{9}





\bibitem{Efr}


Ефремов А.А., {\em Изолированные сверху \re{} степени}, I //


Изв. вузов. Математика. -- 1998. -- \No2. -- С.20--28.


\bibitem{CoYi}


Cooper S.B., Yi X. {\em Isolated d.\,r.\,e. degrees}, University


of Leeds, preprint, 1995.


\bibitem{So}


Soare R.I., {\em Recursively enumerable sets and degrees$:$ a study


of computable functions and computably generated sets}. -- Berlin:


Springer-Verlag, 1987. -- 437 p.





\end{thebibliography}





\vskip 0.5cm





\post{\it Казанский государственный}{\it Поступила}


\post{\it университет}{18.09.1995}





\label{end}


\end{document}


